\section{PRISMA 2020 checklist}
\label{app:prisma}

The review is registered as OSF \texttt{osf.io/ab2wn} (Generalized Systematic Review Registration
v6, approved 2026-09-16, registered 2026-09-16 at 18:12 UTC, before any searching or screening), in
the OSF project \texttt{osf.io/vkjer}. Every departure from the registered text is recorded as a
numbered amendment in \texttt{docs/protocol\_prisma\_p.md} with its date, the protocol section it
changes, and the reason; eleven amendments are accepted at OSF, and a twelfth and a dated revision of it
(12a) are filed and pending. The
amendments are reported in the methods of the main paper (\S4) and logged in full in Supplement~S2,
not here. Nothing in this supplement is a deviation report: this is the reporting map.

Tables~\ref{tab:prisma-title}--\ref{tab:prisma-discussion} give the 27-item PRISMA 2020
checklist~\citep{page2021prisma,page2021prismaexplanation} with the place where each item is reported,
listing the checklist's sub-items separately where they are reported in different places. A section
sign (\S) refers to a section of the main paper; S1--S6 are the sections of this supplement. Every row
was checked against the drafted text of the section it names, and a row whose item the manuscript does
not report says so rather than naming a section. Search reporting additionally follows
PRISMA-S~\citep{rethlefsen2021prismas}, and the protocol was written to
PRISMA-P~\citep{shamseer2015prismap}.

Three items need a word on how they are satisfied, because this review's unit of analysis is a
software system rather than a clinical study. Item~11, risk of bias in the included studies, is
assessed against the five cross-cutting caveats that the main paper's reporting-bias methods (\S4)
name from \texttt{docs/benchmark\_caveats.md} --- benchmark versioning, contamination, self-reporting,
harness--model entanglement, and scaffold drift --- each reported with the direction in which it moves
the estimate and with the analyses it binds. No risk-of-bias instrument is applied to individual
sources and no per-source rating is tabulated. Item~15, certainty, is a \emph{declared absence}:
the same methods subsection states that no formal GRADE-style rating is applied, because the outcome
evidence is author-reported scores on heterogeneous benchmarks rather than comparative trials, and
states what is reported per estimate in its place --- the direction in which the estimate's bias runs,
the discount applied to it, the prediction interval beside the confidence interval, and, for every
null, the minimum detectable effect. Item~22 inherits that decision. Items 25 and 26, support and
competing interests, belong to the front matter of the submitted version and are not yet written; the
table marks them as pending rather than as reported.

Four further items are reported in part, and the tables say which part. The full Boolean search
strings and their per-source translations (item~7) are released in
\texttt{data/raw/search\_log.md} and are not reproduced in the manuscript, which summarizes the
strategy and gives the per-source hit counts. The five full-text exclusion codes and their
sub-reasons (item~16b) are named in the eligibility criteria of \S4 and their counts appear in
the PRISMA flow diagram of \S4 rather than in prose; individual excluded reports are not cited, and
every full-text decision, its deciding step, and the quotes behind it are published in the full-text
audit page (\texttt{screening/fulltext\_audit.html}, whose data is
\texttt{data/screening/fulltext\_audit.js}). Per-system characteristics and per-cell results
(items~17 and~19) are released in full rather than printed: 1,256 systems $\times$ 38 dimensions is
47,728 cells, the manuscript reports them in aggregate and as a 12-system extract (S4), and the
manuscript does not carry a citation for each included system.

\begin{table}[htbp]
\caption{PRISMA 2020 checklist, title, abstract, introduction and methods. \S\ refers to the main
paper; S1--S6 to this supplement.}
\label{tab:prisma-title}
\footnotesize
\begin{tabular}{@{}l p{0.36\linewidth} p{0.40\linewidth}@{}}
\toprule
Item & Checklist item & Reported in \\
\midrule
1 & \raggedright Title: identify the report as a systematic review & \raggedright Title page: ``A Systematic Review, Unified Taxonomy, and Coded Dataset'' \tabularnewline
2 & \raggedright Abstract: structured summary & \raggedright Abstract. One paragraph rather than sub-headed: objectives, methods, results, limitations and the registration identifier, and not support (item~25) \tabularnewline
3 & \raggedright Rationale for the review & \raggedright \S1, \S3 \tabularnewline
4 & \raggedright Objectives: explicit questions the review addresses & \raggedright \S1, RQ1--RQ4 in the registered numbering \tabularnewline
5 & \raggedright Eligibility criteria & \raggedright \S4 (eligibility criteria); the unit of analysis also in \S2 \tabularnewline
6 & \raggedright Information sources, with the date each was last searched & \raggedright \S4 (information sources and search): every source searched 2026-09-16, frozen the same day at 20:59 UTC; the supplementary arm 2026-09-20 \tabularnewline
7 & \raggedright Search strategy: full strings, filters, limits & \raggedright \S4 (information sources and search) in summary, with per-source hit counts. The full strings, per-source translations and filters are in \texttt{data/raw/search\_log.md} and are \emph{not} reproduced in the manuscript \tabularnewline
8 & \raggedright Selection process: who screened, how many, independence, tools & \raggedright \S4 (selection); amendments 3 and 4 in \S4 (registration and amendments) and S2; the models and prompt versions in \S4 (implementation and disclosure) \tabularnewline
9 & \raggedright Data collection process & \raggedright Design cells: \S4 (coding) and S1. Outcome data: \S4 (outcome data), with the ablation harvest described in \S7 (published ablations) \tabularnewline
10a & \raggedright Data items: outcomes sought, and which results were used & \raggedright \S7 (comparable set; published ablations); the released score table is described in S4 \tabularnewline
10b & \raggedright Data items: all other variables sought & \raggedright S1 (all 38 dimensions with their permitted values), \S5 \tabularnewline
11 & \raggedright Study risk of bias assessment & \raggedright \S4 (reporting bias and certainty), \S9 (benchmark caveats): the five cross-cutting caveats, each with its direction. No risk-of-bias instrument and no per-source rating \tabularnewline
12 & \raggedright Effect measures & \raggedright \S4 (synthesis methods); each measure named again where it is used in \S7 \tabularnewline
13a--13d & \raggedright Synthesis methods: eligibility for synthesis, preparation of data, tabulation, model and heterogeneity & \raggedright \S4 (synthesis methods); applied in \S6 and \S7 \tabularnewline
13e--13f & \raggedright Synthesis methods: subgroup analyses, sensitivity analyses & \raggedright \S4 (reporting bias and certainty); applied in \S7 (publication bias) \tabularnewline
14 & \raggedright Reporting bias assessment & \raggedright \S4 (reporting bias and certainty) \tabularnewline
15 & \raggedright Certainty assessment & \raggedright \textbf{Declared absent.} \S4 (reporting bias and certainty) states that no GRADE-style rating is applied, gives the reason, and names what is reported per estimate instead \tabularnewline
\bottomrule
\end{tabular}
\end{table}

\begin{table}[htbp]
\caption{PRISMA 2020 checklist, results. \S\ refers to the main paper; S1--S6 to this supplement.}
\label{tab:prisma-results}
\footnotesize
\begin{tabular}{@{}l p{0.36\linewidth} p{0.40\linewidth}@{}}
\toprule
Item & Checklist item & Reported in \\
\midrule
16a & \raggedright Study selection: records screened, included, excluded, with a flow diagram & \raggedright The PRISMA flow diagram in \S4; the stage counts in \S4 (information sources and search; selection), from \texttt{data/prisma\_counts.json} \tabularnewline
16b & \raggedright Studies that appeared to meet the criteria but were excluded, with reasons & \raggedright The five codes and their sub-reasons in \S4 (eligibility criteria); their counts in the flow diagram and, in one sentence, in \S4 (selection): 1,028 out of scope, 156 no harness description, 82 duplicate systems, 76 not retrievable, 8 other. Individual excluded reports are not cited; every decision with its quotes is in \texttt{screening/fulltext\_audit.html} \tabularnewline
17 & \raggedright Study characteristics and citations & \raggedright In aggregate in \S6, with a 12-system extract in S4 and every system's record in \texttt{data/systems.json}. \emph{Not} a per-system citation list \tabularnewline
18 & \raggedright Risk of bias in the included studies & \raggedright \S9 (benchmark caveats; self-reporting) and \S7 (publication bias): which caveats bind which design, and the measured self-reporting. No per-source ratings \tabularnewline
19 & \raggedright Results of individual studies & \raggedright The 12-system extract in S4; all 47,728 cells in \texttt{data/systems.json} and every extracted score in \texttt{data/results.csv}. \emph{Not} printed per system \tabularnewline
20a & \raggedright Results of syntheses: characteristics of the contributing studies & \raggedright \S7: keys, observations and systems per contrast (comparable set); papers and contrasts per dimension (published ablations) \tabularnewline
20b--20c & \raggedright Results of syntheses: each synthesis, and heterogeneity & \raggedright \S6; \S7 (comparable set; published ablations, with $I^2$ and a prediction interval per pooled dimension; ablation coverage; within-paper comparisons) \tabularnewline
20d & \raggedright Results of syntheses: sensitivity analyses & \raggedright \S7 (published ablations): the discount rule built from trim-and-fill and one-null-per-contrast shrinkage, its Hartung--Knapp variant, and the prediction-interval criterion; \S7 (publication bias): the paper-clustered null of the own-arm test. The per-dimension null-fill counts and the size-based test are in the replication package (S5), not in the manuscript \tabularnewline
21 & \raggedright Reporting biases among the synthesized results & \raggedright \S7 (publication bias); the methods in \S4 (reporting bias and certainty) \tabularnewline
22 & \raggedright Certainty of evidence for each outcome & \raggedright \textbf{No certainty ratings} (item~15). Per estimate instead: bound direction, discount, prediction interval and minimum detectable effect in \S7 (the table of the three designs and their bound directions; the pooled-ablation table; the cross-paper contrasts, reported in prose with their minimum detectable effects) and \S9 \tabularnewline
\bottomrule
\end{tabular}
\end{table}

\begin{table}[htbp]
\caption{PRISMA 2020 checklist, discussion and other information. \S\ refers to the main paper;
S1--S6 to this supplement.}
\label{tab:prisma-discussion}
\footnotesize
\begin{tabular}{@{}l p{0.36\linewidth} p{0.40\linewidth}@{}}
\toprule
Item & Checklist item & Reported in \\
\midrule
23a & \raggedright Interpretation of the results in the context of other evidence & \raggedright \S8; the takeaways of \S6 and \S7; the prior work in \S3 \tabularnewline
23b & \raggedright Limitations of the evidence included & \raggedright \S9 (benchmark caveats; self-reporting), \S7 (ablation coverage) \tabularnewline
23c & \raggedright Limitations of the review processes used & \raggedright \S9: no human coder, reproducibility of the coding, the absence-rule rewrite, the loop-primitives dimension, design weights and frame completeness, construct validity, search bias and the search freeze \tabularnewline
23d & \raggedright Implications for practice, policy and future research & \raggedright \S8 (the reporting checklist; three open questions) \tabularnewline
24a & \raggedright Registration name and number & \raggedright \S4 (registration and amendments): OSF Generalized Systematic Review Registration v6 at \texttt{osf.io/ab2wn}, registered 2026-09-16 at 18:12 UTC. No DOI yet: it is minted when the embargo releases \tabularnewline
24b & \raggedright Where the protocol can be accessed & \raggedright \S4 (registration and amendments): \texttt{osf.io/ab2wn}, embargoed to 2026-10-31 and released on posting of this preprint; \texttt{docs/protocol\_prisma\_p.md} is the working copy \tabularnewline
24c & \raggedright Amendments to the registered information & \raggedright \S4 (registration and amendments) and the full log in S2: eleven numbered amendments accepted at OSF, and a twelfth, covering the controlled ablation of \S7, and a dated revision of it (12a, 2026-09-25), both filed and pending acceptance \tabularnewline
25 & \raggedright Support: sources of financial or non-financial support & \raggedright \textbf{Not yet in the manuscript.} Front matter, to be supplied at submission \tabularnewline
26 & \raggedright Competing interests & \raggedright \textbf{Not yet in the manuscript.} Front matter, to be supplied at submission \tabularnewline
27 & \raggedright Availability of data, code, and other materials & \raggedright The unnumbered ``Data and code availability'' section of the main paper, S4, and the artefact map of S5 \tabularnewline
\bottomrule
\end{tabular}
\end{table}


% TODO: items 25 (support) and 26 (competing interests) are front matter and are not yet in the
% manuscript; add them before submission and change the two rows above.
% NOTE 2026-09-25 (v2): this file is now Supplement S3 and compiles in supplement.tex, a standalone
% build, so pointers into the main paper are written as "\S<n> (topic)" in prose rather than \ref.
% The section numbers follow the v2 outline (1 intro ... 9 threats, 10 conclusion); the topic names
% in parentheses are subsection topics, not titles. If the v2 section order changes, update the \S
% numbers here. Item 27 names the unnumbered "Data and code availability" section by name.
% VERIFIED 2026-09-25: every row above was checked against the drafted text of the section it names
% (v1 manuscript). Rows 7, 9, 10a, 16b, 17, 19 state what the manuscript does and does not report;
% rows 11, 15 and 22 describe declared absences with their rationale; 13, 20 and 23 are split
% because their sub-items are reported in different places. Rows 16b and the item-16b paragraph now
% name screening/fulltext_audit.html, which is where the page is; data/screening/ holds its data file.
% RESOLVED 2026-09-25: data/prisma_counts.json records included_systems = 6,504 and S4's
% Table~\ref{tab:release-counts} carries 6,504 from it.
